% Chapter 1: Euler's Totient Function and Theorems
% ch01_totient.tex

\chapter{Euler's Totient Function and Theorems}

\begin{historical}
Euler introduced the totient function $\varphi(n)$ (also written $\phi(n)$) in 1763,
counting the positive integers $\le n$ that are coprime to $n$. This function
generalizes Fermat's little theorem and plays a central role in modern cryptography.
The function appears in Euler's work on the RSA cryptosystem predecessor and
in the theory of cyclotomic fields.
\end{historical}

\section{Prerequisites}

To work through this chapter, you need:
\begin{itemize}
    \item \textbf{Mathematics:} Basic divisibility and congruence arithmetic, the Euclidean algorithm for $\gcd$, understanding of prime factorization.
    \item \textbf{Programming:} Python functions, loops, and modular exponentiation via the built-in `pow(base, exp, mod)` function.
\end{itemize}

\section{Definition and Basic Properties}

The \textbf{Euler totient function} $\varphi(n)$ counts the integers $1 \le k \le n$ with $\gcd(k,n) = 1$.

\[
\varphi(n) = |\{k \in \mathbb{Z} : 1 \le k \le n,\ \gcd(k,n) = 1\}|
\]

For $n = p_1^{e_1} p_2^{e_2} \cdots p_k^{e_k}$, the formula is:
\[
\varphi(n) = n \prod_{p|n} \left(1 - \frac{1}{p}\right)
\]

\begin{lstlisting}[caption={Computing the totient function}, label=lst:totient]
from euler import euler_totient, euler_totient_sieve, euler_totient_theorem

# Individual values
for n in range(1, 21):
    print(f"\varphi({n:2d}) = {euler_totient(n):2d}")

# Sieve for many values
phi = euler_totient_sieve(30)
print(f"\n\varphi(1..30) = {phi[1:]}")
\end{lstlisting}

\section{Multiplicativity}

\begin{theorem}[Multiplicativity of $\varphi$]
If $\gcd(m,n) = 1$, then $\varphi(mn) = \varphi(m)\varphi(n)$.
\end{theorem}

\begin{proof}
This follows from the Chinese Remainder Theorem. The map $x \mapsto (x \bmod m, x \bmod n)$
gives a bijection between $(\mathbb{Z}/mn\mathbb{Z})^\times$ and $(\mathbb{Z}/m\mathbb{Z})^\times \times (\mathbb{Z}/n\mathbb{Z})^\times$ when $\gcd(m,n) = 1$.
\end{proof}

\begin{theorem}[Totient of Prime Powers]
For a prime $p$ and positive integer $e$:
\[
\varphi(p^e) = p^e - p^{e-1} = p^{e-1}(p-1)
\]
\end{theorem}

\begin{proof}
The integers $1, 2, \ldots, p^e$ that are NOT coprime to $p^e$ are exactly those divisible by $p$.
These are $p, 2p, 3p, \ldots, p^{e-1} \cdot p = p^e$, giving $p^{e-1}$ such integers.
Therefore $\varphi(p^e) = p^e - p^{e-1}$.
\end{proof}

\section{Computational Values}

\begin{table}[H]
\centering
\begin{tabular}{c|c|c|c|c|c|c}
$n$ & 1 & 2 & 3 & 4 & 5 & 6 \\
\hline
$\varphi(n)$ & 1 & 1 & 2 & 2 & 4 & 2 \\
\end{tabular}
\caption{Small values of Euler's totient function}
\label{tab:totient_small}
\end{table}

\begin{lstlisting}[caption={Totient function values}, label=lst:totient_values]
from euler import euler_totient

print("n    \varphi(n)    Reason")
print("---  ------  ---------------------------")
for n in [1, 2, 3, 4, 5, 6, 7, 8, 9, 10]:
    if n == 1:
        reason = "\varphi(1) = 1 by convention"
    elif n == 2:
        reason = "\varphi(2) = 1 (only 1)"
    elif n == 3:
        reason = "\varphi(3) = 2 (1,2)"
    else:
        reason = f"\varphi({n}) computed"
    print(f"{n:3d}  {euler_totient(n):6d}  {reason}")
\end{lstlisting}

\section{Euler's Theorem}

\textbf{Theorem (Euler, 1763).} If $\gcd(a,n) = 1$, then $a^{\varphi(n)} \equiv 1 \pmod{n}$.

This generalizes Fermat's little theorem (where $n=p$ is prime, so $\varphi(p) = p-1$).

\begin{proof}
Let $a$ be coprime to $n$. Consider the set $S = \{r_1, r_2, \ldots, r_{\varphi(n)}\}$ of integers
in $\{1, 2, \ldots, n\}$ coprime to $n$. Multiplying each element by $a$ gives the set
$aS = \{ar_1, ar_2, \ldots, ar_{\varphi(n)}\}$. Since $\gcd(a,n) = 1$, each $ar_i$ is also coprime to $n$.
Moreover, if $ar_i \equiv ar_j \pmod{n}$, then $r_i \equiv r_j \pmod{n}$ (since $a$ is invertible),
so the elements of $aS$ are all distinct modulo $n$.

Therefore $aS$ is a permutation of $S$ modulo $n$, and:
\[
\prod_{i=1}^{\varphi(n)} ar_i \equiv \prod_{i=1}^{\varphi(n)} r_i \pmod{n}
\]
Since each $r_i$ is coprime to $n$, we can cancel $\prod r_i$ from both sides:
\[
a^{\varphi(n)} \equiv 1 \pmod{n}
\]
\end{proof}

\begin{lstlisting}[caption={Euler's theorem and Fermat's little theorem}, label=lst:euler_theorem]
from euler import euler_totient_theorem, fermat_little_theorem, euler_totient

print("Euler's theorem: a^\varphi(n) \equiv 1 (mod n) for gcd(a,n)=1")
for a, n in [(2, 9), (3, 10), (5, 12), (7, 15)]:
    phi = euler_totient(n)
    result = pow(a, phi, n)
    print(f"  {a}^{phi} \equiv {result} (mod {n})  [gcd({a},{n})={__import__('math').gcd(a,n)}]")

print("\nFermat's little theorem: a^(p-1) \equiv 1 (mod p)")
for a, p in [(2, 7), (3, 11), (5, 13), (7, 13)]:
    print(f"  {a}^{p-1} \equiv {pow(a, p-1, p)} (mod {p})")
\end{lstlisting}

\section{The Totient Sum Function}

The sum of totient values over divisors:

\begin{theorem}
$\sum_{d|n} \varphi(d) = n$
\end{theorem}

\begin{proof}
Group the integers $1, 2, \ldots, n$ by their gcd with $n$. For each divisor $d$ of $n$,
the numbers $k$ with $\gcd(k,n) = d$ are those of the form $k = dm$ where $\gcd(m, n/d) = 1$.
There are $\varphi(n/d)$ such values, so:
\[
n = \sum_{k=1}^n 1 = \sum_{d|n} \varphi(d)
\]
\end{proof}

\begin{lstlisting}[caption={Sum of totients over divisors}, label=lst:totient_sum]
from euler import euler_totient

print("Verifying sum_{d|n} \varphi(d) = n:")
for n in [1, 2, 3, 4, 5, 6, 10, 12, 24]:
    divisors = [d for d in range(1, n+1) if n % d == 0]
    sum_phi = sum(euler_totient(d) for d in divisors)
    print(f"  n={n:2d}: divisors={divisors}, sum \varphi = {sum_phi}, n = {n}, match = {sum_phi == n}")
\end{lstlisting}

\section{Applications to Cryptography}

The totient function is fundamental in:
\begin{itemize}
    \item \textbf{RSA cryptography:} Key generation uses $\varphi(n) = (p-1)(q-1)$
    \item \textbf{Primitive roots:} $\varphi(p-1)$ primitive roots modulo prime $p$
    \item \textbf{Cyclotomic polynomials:} $\Phi_n(x)$ has degree $\varphi(n)$
\end{itemize}

\subsection{RSA Key Generation}

The RSA cryptosystem works as follows:

\begin{enumerate}
\item Choose two large primes $p$ and $q$
\item Compute $n = pq$ and $\varphi(n) = (p-1)(q-1)$
\item Choose $e$ such that $\gcd(e, \varphi(n)) = 1$
\item Compute $d = e^{-1} \pmod{\varphi(n)}$
\item Public key: $(n, e)$, Private key: $(n, d)$
\end{enumerate}

Encryption: $c \equiv m^e \pmod{n}$\
Decryption: $m \equiv c^d \pmod{n}$

\begin{theorem}[RSA Correctness]
For any message $m < n$, decrypting the encrypted message recovers $m$:
\[
(c^d \bmod n) = (m^e)^d \bmod n = m^{ed} \bmod n = m
\]
since $ed \equiv 1 \pmod{\varphi(n)}$ implies $ed = 1 + k\varphi(n)$ for some integer $k$,
and by Euler's theorem $m^{k\varphi(n)} \equiv 1 \pmod{n}$ when $\gcd(m,n) = 1$.
\end{theorem}

\begin{lstlisting}[caption={RSA key generation}, label=lst:rsa]
from euler import euler_totient
import random

# Generate small RSA keys for demonstration
def generate_rsa_keys(bits=16):
    # In practice, use cryptographically secure primes
    p = 65537  # A common choice for p
    q = 65521
    n = p * q
    phi = (p - 1) * (q - 1)
    
    # Choose e (commonly 65537)
    e = 65537
    while euler_totient(e) != euler_totient(e):
        e = random.randrange(3, phi, 2)
        if __import__('math').gcd(e, phi) == 1:
            break
    
    # Compute d = e^(-1) mod phi
    d = pow(e, -1, phi)
    return n, e, d

n, e, d = generate_rsa_keys()
print(f"RSA modulus n = {n}")
print(f"Public exponent e = {e}")
print(f"Private exponent d = {d}")
print(f"Verification: e*d mod \varphi(n) = {e*d % 5948838576} (should be 1)")
\end{lstlisting}

\section{Modern Relevance}

The totient function is not merely a historical curiosity---it is the backbone of modern public-key cryptography. The RSA encryption scheme, used to secure virtually every online transaction, email, and government communication, relies on the computational difficulty of computing $\varphi(n)$ without knowing the prime factorization of $n$. Given a public key $(n, e)$, an attacker who can efficiently compute $\varphi(n)$ can recover the private key and break the encryption. This connection between an 18th-century number-theoretic function and 21st-century cybersecurity is one of the most striking examples of pure mathematics serving practical needs.

\section{Computational Complexity}

\subsection{Computing $\varphi(n)$}
\begin{itemize}
\item Naive method: $O(\sqrt{n})$ (trial division)
\item Sieve method: $O(n \log \log n)$ for all values up to $n$
\item Using prime factorization: $O(\sqrt{n})$ per value, but fast after factorization
\end{itemize}

\subsection{RSA Key Generation}
\begin{itemize}
\item Time to generate key: $O(\log^3 n)$ for $n$-bit modulus (probabilistic primality testing)
\item Security: Based on hardness of factoring $n = pq$ and computing $\varphi(n)$
\end{itemize}

\section{Frequently Asked Questions}

\begin{itemize}
\item \textbf{Q: Why is $\varphi(1) = 1$?} \\
The integer $1$ is coprime to itself ($\gcd(1,1) = 1$), so it is counted. Thus $\varphi(1) = 1$.

\item \textbf{Q: Is the totient function multiplicative?} \\
Yes, but only \emph{completely} for coprime arguments: if $\gcd(m,n) = 1$, then $\varphi(mn) = \varphi(m)\varphi(n)$. If $\gcd(m,n) \neq 1$, this equality generally fails. For example, $\varphi(2 \cdot 3) = \varphi(6) = 2 = \varphi(2)\varphi(3)$, but $\varphi(2 \cdot 4) = \varphi(8) = 4 \neq \varphi(2)\varphi(4) = 1 \cdot 2 = 2$.

\item \textbf{Q: How do I compute $\varphi(n)$ for a large $n$?} \\
Factor $n$ into primes: $n = p_1^{e_1} \cdots p_k^{e_k}$. Then $\varphi(n) = n \prod_{i=1}^k (1 - 1/p_i)$. The bottleneck is factoring $n$, not the formula itself. For many values, a sieve (as shown in the code) is more efficient.

\item \textbf{Q: Why does Euler's theorem generalize Fermat's little theorem?} \\
Fermat's little theorem is the special case $n = p$ (prime), where $\varphi(p) = p-1$. Euler's theorem holds for any $n$ with $\gcd(a,n) = 1$, reducing to Fermat's when $n$ is prime.

\item \textbf{Q: What is the maximum value of $\varphi(n)/n$?} \\
$\varphi(n)/n = \prod_{p|n}(1 - 1/p)$, which is maximized when $n$ is prime, giving $(p-1)/p \to 1$ as $p \to \infty$. The minimum occurs when $n$ has many small prime factors; for example, $\varphi(n)/n \geq 1/\log\log n$ (Robin's inequality).

\item \textbf{Q: How is the totient function used in RSA?} \\
In RSA, $n = pq$ for large primes $p, q$. The public key exponent $e$ and private key $d$ satisfy $ed \equiv 1 \pmod{\varphi(n)}$, where $\varphi(n) = (p-1)(q-1)$. Encryption is $c \equiv m^e \pmod{n}$ and decryption is $m \equiv c^d \pmod{n}$.

\item \textbf{Q: What is Carmichael's function?} \\
Carmichael's function $\lambda(n)$ is the smallest positive integer such that $a^{\lambda(n)} \equiv 1 \pmod{n}$ for all $\gcd(a,n) = 1$. It divides $\varphi(n)$ and is equal to $\varphi(n)$ for prime powers $p^e$ where $p$ is odd or $e \leq 2$. For $n = 2^e$ with $e \geq 3$, $\lambda(n) = 2^{e-2}$.

\item \textbf{Q: Can I use the totient function to find primitive roots?} \\
A number $g$ is a primitive root modulo $p$ if the order of $g$ modulo $p$ is $p-1 = \varphi(p)$. There are $\varphi(\varphi(p)) = \varphi(p-1)$ primitive roots modulo $p$. To find one, test candidates by checking that $g^{\varphi(p)/q} \not\equiv 1 \pmod{p}$ for each prime factor $q$ of $\varphi(p)$.
\end{itemize}

\section{Exercises}

\begin{exercise}
Prove that $\sum_{d|n} \varphi(d) = n$ by counting the elements of $\mathbb{Z}/n\mathbb{Z}$ according to their order.
\end{exercise}

\begin{exercise}
Show that the probability two random integers are coprime is $6/\pi^2 = 1/\zeta(2)$.
\end{exercise}

\begin{exercise}
Implement the Chinese Remainder Theorem using the totient function.
\end{exercise}

\begin{exercise}[Computational]
Optimize the totient sieve function \texttt{euler\_totient\_sieve} to run in linear time $O(N)$ using a prime sieve.
\end{exercise}

\begin{exercise}
Prove that $\varphi(n)$ is even for all $n \geq 3$.
\end{exercise}

\begin{exercise}
Show that $\varphi(n) \geq \sqrt{n/2}$ for all $n \geq 1$.
\end{exercise}

\begin{exercise}
Prove that $\varphi(n) = n \prod_{p|n} (1 - 1/p)$ using the multiplicativity property and the formula for prime powers.
\end{exercise}

\begin{exercise}[Computational]
Implement a function to compute $\varphi(n)$ using the prime factorization of $n$. Compare its performance with the sieve approach for $n$ up to $10^6$.
\end{exercise}

\begin{exercise}
Show that the number of solutions to $x^k \equiv 1 \pmod{p}$ is $\gcd(k, p-1)$ for prime $p$.
\end{exercise}

\begin{exercise}
Prove that if $p$ is an odd prime, then every primitive root modulo $p$ is a quadratic non-residue.
\end{exercise}

\section{Totient Applications in Cryptography}

\subsection{The RSA Cryptosystem}

The security of the RSA cryptosystem rests on the computational difficulty of factoring large integers and computing $\varphi(n)$ without knowing the prime factorization. The protocol works as follows:

\begin{enumerate}
    \item \textbf{Key Generation:} Choose two large primes $p$ and $q$, compute $n = pq$ and $\varphi(n) = (p-1)(q-1)$.
    \item \textbf{Public Key:} Select $e$ such that $\gcd(e, \varphi(n)) = 1$, typically $e = 65537$.
    \item \textbf{Private Key:} Compute $d = e^{-1} \pmod{\varphi(n)}$.
\end{enumerate}

Encryption: $c \equiv m^e \pmod{n}$\\
Decryption: $m \equiv c^d \pmod{n}$

The correctness follows from Euler's theorem: since $ed \equiv 1 \pmod{\varphi(n)}$, we have $ed = 1 + k\varphi(n)$ for some integer $k$, and $m^{ed} \equiv m \pmod{n}$ when $\gcd(m,n) = 1$.

\begin{lstlisting}[caption={RSA encryption and decryption}, label=lst:rsa_crypto]
from euler import euler_totient

# Small RSA for demonstration
p, q = 61, 53
n = p * q
phi = (p - 1) * (q - 1)  # Euler's totient

e = 17  # Public exponent
d = pow(e, -1, phi)  # Private exponent

# Encrypt
m = 42
c = pow(m, e, n)
print(f"Encrypted: {c}")

# Decrypt
m_recovered = pow(c, d, n)
print(f"Decrypted: {m_recovered}")
\end{lstlisting}

\subsection{Diffie-Hellman Key Exchange}

The Diffie-Hellman key exchange allows two parties to establish a shared secret over an insecure channel. The security relies on the discrete logarithm problem, which is related to the structure of $\mathbb{Z}_p^*$ whose order is $\varphi(p) = p-1$.

\begin{lstlisting}[caption={Diffie-Hellman key exchange}, label=lst:dh]
p = 23  # Prime modulus
g = 5  # Generator

# Alice's secret and public key
a = 6
A = pow(g, a, p)

# Bob's secret and public key
b = 7
B = pow(g, b, p)

# Shared secret
s_alice = pow(B, a, p)
s_bob = pow(A, b, p)
print(f"Shared secret: {s_alice}")  # Both get the same value
\end{lstlisting}

\subsection{Carmichael Function}

Carmichael's function $\lambda(n)$ is the smallest positive integer such that $a^{\lambda(n)} \equiv 1 \pmod{n}$ for all $\gcd(a,n) = 1$. It satisfies $\lambda(n) | \varphi(n)$, and for RSA key generation, $\lambda(n) = \text{lcm}(p-1, q-1)$ when $n = pq$.

\begin{theorem}[Carmichael's Theorem]
For $n = p_1^{e_1} \cdots p_k^{e_k}$, we have:
\[
\lambda(n) = \text{lcm}(\lambda(p_1^{e_1}), \ldots, \lambda(p_k^{e_k}))
\]
where $\lambda(p^e) = \varphi(p^e)$ for odd primes, and $\lambda(2) = 1$, $\lambda(4) = 2$, $\lambda(2^e) = 2^{e-2}$ for $e \\geq 3$.
\end{theorem}

\section{Totient in Number Theory}

\subsection{Gauss's Circle Problem}

The totient function appears in counting lattice points visible from the origin. The number of visible lattice points in the first quadrant within a circle of radius $R$ is asymptotically $\frac{3}{\pi^2} R^2$.

\subsection{Reduced Residue Systems}

A \textbf{reduced residue system modulo $n$} is a set of $\varphi(n)$ integers that are each coprime to $n$ and pairwise incongruent modulo $n$. For example, $\\{1, 3, 7, 9\\}$ is a reduced residue system modulo $10$.

\begin{theorem}
The product of all elements in a reduced residue system modulo $n$ is congruent to:
\[
\begin{cases}
-1 \pmod{n} & \text{if } n = 4, p^k, \text{ or } 2p^k \text{ for odd prime } p \\
1 \pmod{n} & \text{otherwise}
\end{cases}
\]
This generalizes Wilson's theorem.
\end{theorem}

\subsection{Multiplicative Order}

The \textbf{multiplicative order} of $a$ modulo $n$ is the smallest positive integer $k$ such that $a^k \equiv 1 \pmod{n}$. By Lagrange's theorem applied to the group $(\mathbb{Z}/n\mathbb{Z})^\times$, the order always divides $\varphi(n)$.

\begin{lstlisting}[caption={Computing multiplicative order}, label=lst:order]
from euler import euler_totient
import math

def multiplicative_order(a, n):
    """Find smallest k such that a^k \equiv 1 (mod n)"""
    if math.gcd(a, n) != 1:
        return None
    phi = euler_totient(n)
    # Check divisors of \varphi(n)
    for d in sorted([k for k in range(1, phi + 1) if phi % k == 0]):
        if pow(a, d, n) == 1:
            return d
    return phi

print(f"Order of 2 mod 7: {multiplicative_order(2, 7)}")  # 3
print(f"Order of 3 mod 7: {multiplicative_order(3, 7)}")  # 6 (primitive root)
\end{lstlisting}

\section{Sums and Averages}

\subsection{Asymptotic Behavior}

The average order of $\varphi(n)$ satisfies:
\[
\sum_{n \leq x} \varphi(n) = \frac{3x^2}{\pi^2} + O(x \\log x)
\]

This follows from the identity $\sum_{d|n} \varphi(d) = n$ and properties of the Riemann zeta function.

\subsection{Mean Value}

The ratio $\varphi(n)/n$ has average value $6/\pi^2$ over large ranges, reflecting the probability that two random integers are coprime.

\begin{lstlisting}[caption={Average totient values}, label=lst:avg_totient]
from euler import euler_totient_sieve
import math

N = 1000
phi = euler_totient_sieve(N)
avg_ratio = sum(phi[n]/n for n in range(1, N+1)) / N
print(f"Average \varphi(n)/n for n \leq {N}: {avg_ratio:.4f}")
print(f"Theoretical value 6/\pi^2: {6/math.pi**2:.4f}")
\end{lstlisting}

\section{Exercises}

\begin{exercise}
Prove Gauss's generalization of Wilson's theorem: the product of elements in a reduced residue system modulo $n$ is $-1$ if $n = 4, p^k$, or $2p^k$, and $1$ otherwise.
\end{exercise}

\begin{exercise}
Show that the multiplicative order of $a$ modulo $n$ divides $\varphi(n)$ for all $\gcd(a,n) = 1$.
\end{exercise}

\begin{exercise}
Compute $\lambda(n)$ for $n = 15, 24, 35$ and verify that $\lambda(n) | \varphi(n)$.
\end{exercise}

\begin{exercise}[Computational]
Implement a function to compute the Carmichael function $\lambda(n)$ using the prime factorization. Test it for $n \leq 100$.
\end{exercise}

\begin{exercise}
Prove that if $p$ is an odd prime, then the number of primitive roots modulo $p$ is $\varphi(p-1)$.
\end{exercise}

\begin{exercise}
Show that $\varphi(n) \leq n - \\sqrt{n}$ for all composite $n > 1$.
\end{exercise}

\begin{exercise}
Compute the sum $\sum_{n=1}^{100} \varphi(n)$ and verify the asymptotic formula $\frac{3 \\cdot 100^2}{\pi^2} \\approx 303.96$.
\end{exercise}

\begin{exercise}
Prove that the number of irreducible polynomials of degree $n$ over $\mathbb{F}_p$ is $\frac{1}{n} \sum_{d|n} \\mu(d) p^{n/d}$, where $\\mu$ is the Moebius function.
\end{exercise}
